\documentclass{situationist-preprint}

\title{Buddhist Reductionism and Parfit}
\msid{PID-13-b}
\msfamily{PID-13}
\msdiscipline{Personal Identity and the Self}
\mstype{Companion manuscript}
\mskeywords{no-self; Buddhist reductionism; Parfit; two truths; process view; impermanence}

\begin{document}
\maketitle

\begin{abstract}
PID-13-a argues that there is no permanent self: the self is a process that persists only as long as its constituent activities. Two traditions reached a similar conclusion by different routes: the Buddhist doctrine of no-self, in its Abhidharma form, and Parfit's reductionism. We compare them and locate the view of PID-13-a between them. All three deny a further fact beyond physical and psychological events and their relations. The Buddhist reductionist adds that persons exist only conventionally; Parfit adds that identity is not what matters. PID-13-a accepts Parfit's addition and rejects the Buddhist one: the self is a process, and a process exists as fully as anything in time does.
\end{abstract}

\section{Two reductionisms}

In the \emph{Milindapa\~nha}, the monk N\=agasena answers King Milinda's question about his identity by asking whether the chariot is the pole, the axle, the wheels or the frame, or something besides them. It is none of these; \enquote{chariot} is a designation used for the parts suitably assembled, and \enquote{N\=agasena} is a designation used for the aggregates of body, feeling, perception, volitional formations and consciousness \citep{RhysDavids1890}. The Abhidharma systematized the point into a doctrine of two truths: ultimately there exist only momentary psycho-physical events, the \emph{dharmas}; persons, like chariots, exist only conventionally, as useful designations for causally connected series of such events \citep{Siderits2003,Collins1982}.

\citet[ch.~11 and App.~J]{Parfit1984}, who noted the parallel himself, defended a reductionist view on which a person's existence consists in the existence of a brain and body and the occurrence of a series of interrelated physical and mental events, and on which a complete description of reality could be given in impersonal terms. He drew a practical conclusion: since there is no further fact, identity is not what matters, and the belief that it is has made us care too much about our own continued existence.

\begin{proposition}\label{p:common}
Buddhist reductionism, Parfit's reductionism and PID-13-a agree that no entity distinct from the organism and the integrating activities plays the self-role and persists through every change in, and cessation of, those activities.
\end{proposition}

\begin{proof}
For the Buddhist reductionist, the ultimate constituents are momentary events, and nothing among them persists. For Parfit, there is no separately existing entity, distinct from brain, body and events, and no further fact. PID-13-a, Lemma~3.2, holds that no bearer other than the organism plays the self-role, and Lemma~3.3 that the organism is not a permanent self. So none of the three admits a permanent self (PID-13-a, Definition~2.2).
\end{proof}

\section{Where they diverge}

\paragraph{Conventional existence.} The Buddhist reductionist holds that the person, and the self, exist only by convention: the designation is useful, but nothing ultimately answers to it. Parfit at times wrote as if persons were real but not \enquote{separately existing}; he did not say that they exist only by convention. PID-13-a holds that the self is a process, and PID-04 that it is a construction and real. A process, the integrating activity by which a subject binds its experience, retains its past and carries out its intentions (PID-13-a, Assumption~2.4), is not a designation of a series. It figures in explanations that cannot be given in terms of the momentary events alone, such as why a deliberation issues in an action or why a perception is recognized as familiar. On the explanatory criterion of PID-04, it therefore exists. The Buddhist two-truths doctrine and PID-13-a agree about what there is at the level of events; they disagree whether the organized activity of those events is a further real thing, and the explanatory criterion decides in favor of reality.

\paragraph{What matters.} Parfit concluded that identity is not what matters, and PID-03 and PID-06 agree: connectedness carries most of what matters, with a small residual weight for identity. The Buddhist conclusion is more radical: since the self is empty, concern for one's own future as one's own is a form of clinging, to be extinguished. Parfit took a moderate version of this lesson, writing that after he came to his view he cared less about his own death and more about the lives of others. PID-13-a takes no position on the soteriological claim. It is consistent with it that concern for a future being should be proportioned to the connections that carry what matters, and not to an identity that would require a permanent self.

\paragraph{Presence in each act.} A further difference concerns the first-person perspective. The Abhidharma analysis treats the sense of being a subject as itself a constructed event among others. PID-13-a, with PHE-20, recognizes in each act of reflection a subject that the act cannot bracket, present in that act and not before or after it. This is not a permanent self, since it does not outlast the act; but it is not a designation either. It is the integrating process as it is lived in each of its phases. Parfit's impersonal description, which omits it, is in this respect incomplete: an impersonal description can say that an act occurred and was integrated with others, but not, without loss, that it was lived by the one who performs it.

\begin{objection}
If the self is a process and processes exist, the view reinstates the self that Buddhist and Parfitian reductionists deny.
\end{objection}

\begin{reply}
What they deny is a self that is a further fact or a permanent entity. A process is neither. It consists in the activities that constitute it and exists only while they go on (PID-13-a, Definition~2.3). To affirm that it exists is to affirm that the integrated activity is real, which no reductionist denies at the level of events; it is not to affirm a bearer behind the activity.
\end{reply}

\section{Conclusion}

Buddhist reductionism and Parfit's reductionism share the denial of a permanent self. PID-13-a shares it too. It stands closer to Parfit than to the Buddhist on the reality of persons, holding that the self is a real process and not merely a designation, and with both in denying that a persisting self is what one's concern should track.

\begin{thebibliography}{99}\small
\bibitem[Collins(1982)]{Collins1982} Collins, S. (1982). \emph{Selfless Persons: Imagery and Thought in Therav\=ada Buddhism}. Cambridge University Press.
\bibitem[Parfit(1984)]{Parfit1984} Parfit, D. (1984). \emph{Reasons and Persons}. Oxford University Press.
\bibitem[Rhys Davids(1890)]{RhysDavids1890} Rhys Davids, T.~W. (trans.) (1890). \emph{The Questions of King Milinda}, Part I. Sacred Books of the East, vol.~35. Clarendon Press.
\bibitem[Siderits(2003)]{Siderits2003} Siderits, M. (2003). \emph{Personal Identity and Buddhist Philosophy: Empty Persons}. Ashgate.
\end{thebibliography}

\end{document}
